\subsection{Centralizers and Normalizers, Stabilizers and Kernels}

\begin{exercise}[2.2.1]
    Prove that $C_G(A)=\set{g\in G\mid g^{-1}ag=a\text{ for all }a\in A}.$
\end{exercise}

\begin{sol}
    $g \in C_G(A) \iff ga = ag$ for all $a \in A \iff g^{-1}ga = g^{-1}ag \iff a = g^{-1}ag$ for all $a \in A$.
\end{sol}

\begin{exercise}[2.2.2]
    Prove that $C_G(Z(G))=G$ and deduce that $N_G(Z(G))=G.$
\end{exercise}

\begin{sol}
    By definition, $C_G(Z(G)) \subseteq G$. Now let $g \in G$ and $z \in Z(G)$. Then $gz = zg$, or $gzg\inv = z$. Then $g \in C_G(Z(G))$, hence $G \subseteq C_G(Z(G))$. It follows that $C_G(Z(G)) = G$. Since $C_G(Z(G)) \leq N_G(Z(G)) \leq G$, then $N_G(Z(G)) = G$.
\end{sol}

\begin{exercise}[2.2.3]
    Prove that if $A$ and $B$ are subsets of $G$ with $A \subseteq B$ then $C_G(B)$ is a subgroup of $C_G(A)$.
\end{exercise}

\begin{sol}
    Pick $g \in C_G(B)$. Then $gbg\inv = b$ for all $b \in B$. In particular, $gag\inv = a$ for all $a \in A$ because $A \subseteq B$. Then $C_G(B) \subseteq C_G(A)$. Since $C_G(B)$ and $C_G(A)$ are subgroups of $G$, then $C_G(B) \leq C_G(A)$.
\end{sol}

\begin{spex}[2.2.4]
    For each of $S_3$, $D_8$, and $Q_8$ compute the centralizers of each element and find the center of each group. Does \defref{Lagrange's theorem} simplify your work?
\end{spex}

\begin{sol}
    Note that $Z(G)$ is the intersection of all centralizers of $G$, so that we can compute $Z(G)$ after computing all centralizers. Note that $C_G(1) = G$ for any group $G$, so we only need to compute the centralizers of the non-identity elements. Beginning with $S_3$, we know by \exref{1.3.20} that $(1\ 2)$ and $(1\ 3)$ generate $S_3$, hence $C_{S_3}((1\ 2))$ cannot be the whole group. Since $\set{1, (1\ 2)} \leq C_{S_3}((1\ 2))$, then $|C_{S_3}((1\ 2))| = 2$. By symmetry, the same argument holds for $(1\ 3)$ and $(2\ 3)$, so that
    \[
    	C_{S_3}((1\ 3)) = \set{1, (1\ 3)}, \quad C_{S_3}((2\ 3)) = \set{1, (2\ 3)}.
    \]
    A similar argument for $(1\ 2\ 3)$ along with the observation that $(1\ 2)(1\ 2\ 3) \neq (1\ 2\ 3)(1\ 2)$ shows that $C_{S_3}((1\ 2\ 3)) = \set{1, (1\ 2\ 3), (1\ 3\ 2)}$. By symmetry, the same argument holds for $(1\ 3\ 2)$. Then $C_{S_3}((1\ 3\ 2)) = \set{1, (1\ 2\ 3), (1\ 3\ 2)}$. Finally, $Z(S_3) = \set 1$.
    
    The following are the centralizers for each element in $D_8$:
    \begin{align*}
        C_{D_8}(1) & = D_8 & C_{D_8}(r) & = \set{1, r, r^2, r^3} \\
        C_{D_8}(r^3) & = \set{1, r, r^2, r^3} & C_{D_8}(r^2) & = D_8 \\
        C_{D_8}(sr) & = \set{1, r^2, sr, sr^3} & C_{D_8}(s) & = \set{1, r^2, s, sr^2} \\
        C_{D_8}(sr^3) & = \set{1, r^2, sr, sr^3} & C_{D_8}(sr^2) & = \set{1, r^2, s, sr^2}
    \end{align*}
    and $Z(D_8) = \set{1, r^2}$. Lastly, the centralizers of $Q_8$ are
    \begin{align*}
        C_{Q_8}(1) & = Q_8 & C_{Q_8}(-1) & = Q_8 \\
        C_{Q_8}(i) & = \set{1, -1, i, -i} & C_{Q_8}(-i) & = \set{1, -1, i, -i} \\
        C_{Q_8}(j) & = \set{1, -1, j, -j} & C_{Q_8}(-j) & = \set{1, -1, j, -j} \\
        C_{Q_8}(k) & = \set{1, -1, k, -k} & C_{Q_8}(-k) & = \set{1, -1, k, -k}
    \end{align*}
    and $Z(Q_8) = \set{1, -1}$.

    Lagrange's theorem simplifies the work by allowing us to determine the possible sizes of the centralizers, which assists in narrowing down candidacy for the centralizers.
\end{sol}

\begin{exercise}[2.2.5]
    In each of parts (a) to (c) show that for the specified group $G$ and subgroup $A$ of $G$, $C_G(A)=A$ and $N_G(A)=G.$
    \begin{subexercises}
        \item $G=S_3$ and $A=\set{1,(1\ 2\ 3),(1\ 3\ 2)}.$
        \item $G=D_8$ and $A=\set{1, s, r^2, sr^2}.$
        \item $G=D_{10}$ and $A=\set{1,r,r^2,r^3,r^4}.$
    \end{subexercises}
\end{exercise}

\begin{solenum}
    \item Observe that $A$ is generated by $(1\ 2\ 3)$. From \exref{2.2.4}, we know $C_{S_3}(A) = A$. Then 3 divides $|N_G(A)|$ and $|N_G(A)|$ divides 6 by Lagrange's Theorem. Consider $(1\ 2) \in S_3$. Then
    \[
    	(1\ 2)A(1\ 2) = \set{(1\ 2)(1)(1\ 2), (1\ 2)(1\ 2\ 3)(1\ 2), (1\ 2)(1\ 3\ 2)(1\ 2)} = \set{1, (1\ 3\ 2), (1\ 2\ 3)} = A
    \]
    so that $(1\ 2) \in N_G(A)$. Hence, $|N_G(A)| = 6$ and $N_G(A) = G$.
    \item Since $|A| = 4$, then 4 divides $|C_G(A)|$ and $|C_G(A)|$ divides 8 by Lagrange's Theorem. Observe that $r \notin C_G(A)$ since $rs = sr^3 \neq sr$. Then $|C_G(A)| \neq 8$, hence $|C_G(A)| = 4$ and $C_G(A) = A$. Then 4 divides $|N_G(A)|$ and $|N_G(A)|$ divides 8 by Lagrange's Theorem. Consider $r \in D_8$. Then
    \[
    	rA r\inv = \set{r(1)r\inv, r(s)r\inv, r(r^2)r\inv, r(sr^2)r\inv} = \set{1, sr^2, r^2, s} = A
    \]
    so that $r \in N_G(A)$. Then $|N_G(A)| = 8$ and $N_G(A) = G$.
    \item Note 5 divides $|C_G(A)|$ and $|C_G(A)|$ divides 10 by Lagrange's Theorem. Since $s \notin C_G(A)$ because $sr \neq rs$, then $C_G(A) = A$. Then 5 divides $|N_G(A)|$ which divides 10 by Lagrange's. A similar computation shows $sAs\inv = \set{1, r^4, r^3, r^2, r} = A$, hence $N_G(A) = G$.
\end{solenum}

\begin{exercise}[2.2.6]
    Let $H$ be a subgroup of the group $G.$
    \begin{subexercises}
        \item Show that $H\le N_G(H).$ Give an example to show that this is not necessarily true if $H$ is not a subgroup.
        \item Show that $H\le C_G(H)$ if and only if $H$ is Abelian.
    \end{subexercises}
\end{exercise}

\begin{solenum}
    \item Let $h, k \in H$. Then $hkh\inv \in H$ by closure. Conversely, let $y \in H$. Then $y = h(h\inv y h)h\inv \in hHh\inv$.
    
    As a counter example, consider the set $G = D_6$ and $H = \set{1, r, s}$. Then
    \[
    	rHr\inv = \set{1, r, sr} \neq H
    \]
    so that $r$ does not normalize $H$, and $H$ is not a subgroup of $N_G(H)$.
    \item 
    \begin{itemize}
        \item [\rightimp] Suppose $H \leq C_G(H)$. Then for any $h_1, h_2 \in H$, we have $h_1 h_2 = h_2 h_1$ since $h_1 \in C_G(H)$. Hence, $H$ is Abelian.
        \item [\leftimp] Suppose $H$ is Abelian, and let $h_1 \in H$. Then for any $h_2 \in H$, we have $h_1 h_2 = h_2 h_1$, so that $h_1 \in C_G(H)$. Since $h_1$ was arbitrary, then $H \leq C_G(H)$. \qh
    \end{itemize}
\end{solenum}

\begin{exercise}[2.2.7] 
    Let $n\in\z$ with $n\ge 3.$ Prove the following:
    \begin{subexercises}
        \item $Z(D_{2n})=\set{1}$ if $n$ is odd.
        \item $Z(D_{2n})=\set{1,r^k}$ if $n=2k.$
    \end{subexercises}
\end{exercise}

\begin{sol}
    We claim that $Z(D_{2n})$ contains no reflections and contains $r^i$ if and only if $n \mid 2i$. To see the first part, suppose $sr^j \in Z(D_{2n})$ for some $0 \leq j < n$. Since it must commute with $r$, then
    \[
    	sr^j r = r sr^j \implies sr^{j + 1} = sr^{j - 1} \implies r^2 = 1.
    \]
    which is a contradiction since $n \geq 3$. Hence, no reflections are in $Z(D_{2n})$.

    Now consider $r^i \in Z(D_{2n})$ for some $0 \leq i < n$. Since it must commute with $s$, then
    \[
    	r^i s = s r^i \implies r^is = r^{-i} s \implies r^{2i} = 1,
    \]
    so that $n \mid 2i$.
    \begin{subexercises}
        \item If $n$ is odd, then $n \mid 2i$ implies that $n \mid i$. Since $0 \leq i < n$, then $i = 0$ and $Z(D_{2n}) = \set 1$.
        \item If $n = 2k$, then $n \mid 2i$ implies that $2k \mid 2i$. Simplifying, we have $k \mid i$. Since $0 \leq i < 2k$, then $i = 0$ or $i = k$, so that $Z(D_{2n}) = \set{1, r^k}$. \qh
    \end{subexercises}
\end{sol}

\begin{exercise}[2.2.8]
    Let $G=S_n$, fix an $i \in \set{1, 2, \dots, n}$ and let $G_i=\set{\sigma\in G\mid\sigma(i)=i}$ (the stabilizer of $i$ in $G$). Use group actions to prove that $G_i$ is a subgroup of $G$. Find $|G_i|$.
\end{exercise}

\begin{sol}
    The fact that $G_i$ stabilizes $i$ under the action of $G$ on $\set{1, 2, \ldots, n}$ by permutation implies that $G_i$ is a subgroup of $G$.

    To find $|G_i|$, observe that any $\sigma \in G_i$ is completely determined by its action on the set $\set{1, 2, \ldots, n} \setminus \set{i}$, which has $n - 1$ elements. Since $\sigma$ is a permutation, then there are $(n - 1)!$ ways to permute these $n - 1$ elements. It follows that $|G_i| = (n - 1)!$.
\end{sol}

\begin{exercise}[2.2.9]
    For any subgroup $H$ of $G$ and any nonempty subset $A$ of $G$ define $N_H(A)$ to be the set $\set{h \in H \mid hAh^{-1} = A}$. Show that $N_H(A) = N_G(A) \cap H$ and deduce that $N_H(A)$ is a subgroup of $H$ (note that $A$ need not be a subset of $H$).
\end{exercise}

\begin{sol}
    Suppose $h \in N_H(A)$. By definition, $hAh\inv = A$ and $h \in H$. Then $h \in N_G(A)$ and $h \in H$, so that $h \in N_G(A) \cap H$. It follows that $N_H(A) \subseteq N_G(A) \cap H$.

    Now suppose $h \in N_G(A) \cap H$. Then $hAh\inv = A$ and $h \in H$. By definition, then $h \in N_H(A)$. It follows that $N_G(A) \cap H \subseteq N_H(A)$. Hence, $N_H(A) = N_G(A) \cap H$. Since $N_G(A) \leq G$ and $H \leq G$, and $N_H(A)$ is the intersection of these two subgroups, then $N_H(A) \leq H$ by \exref{2.1.10}.
\end{sol}

\begin{exercise}[2.2.10]
    Let $H$ be a subgroup of order $2$ in $G$. Show that $N_G(H) = C_G(H)$. Deduce that if $N_G(H) = G$ then $H \le Z(G)$.
\end{exercise}

\begin{sol}
    By assumption, $H = \set{1, h}$ for some $h \in G$ with $h \neq 1$. Since $C_G(H) \leq N_G(H)$, it suffices to show that $N_G(H) \subseteq C_G(H)$. To that end, pick $g \in N_G(H)$. Then $gHg\inv = H$, so that
    \[
    	gHg\inv = \set{g(1)g\inv, g h g\inv} = \set{1, g h g\inv} = H = \set{1, h}
    \]
    It follows that $g h g\inv = h$, so that $g \in C_G(H)$. Since $g$ was arbitrary, then $N_G(H) \subseteq C_G(H)$, hence $N_G(H) = C_G(H)$.

    Now suppose $N_G(H) = G$. Then $C_G(H) = G$, so that for any $g \in G$, we have $g h g\inv = h$. It follows that $h g = g h$ for all $g \in G$, hence $H \leq Z(G)$.
\end{sol}

\begin{exercise}[2.2.11]
    Prove that $Z(G)\le N_G(A)$ for any subset $A$ of $G$.
\end{exercise}

\begin{sol}
    Let $z \in Z(G)$. For any $zaz\inv \in zAz\inv$, then $z \in Z(G)$ implies $zaz\inv = a \in A$. Hence $zAz\inv \subseteq A$. Conversely, for any $a \in A$, we have $a = zaz\inv$ for some $a \in A$ since $z \in Z(G)$. Hence $A \subseteq zAz\inv$. It follows that $zAz\inv = A$, so that $z \in N_G(A)$. Since $z$ was arbitrary, then $Z(G) \leq N_G(A)$.
\end{sol}

\begin{exercise}[2.2.12]
    Let $R$ be the set of all polynomials with integer coefficients in the independent variables $x_1, x_2, x_3, x_4$ i.e., the members of $R$ are finite sums of elements of the form $ax_1^{r_1}x_2^{r_2}x_3^{r_3}x_4^{r_4}$ where $a$ is any integer and $r_1, \ldots, r_4$ are nonnegative integers. For example,
    \begin{equation}
        \label{ex2.2.11eq}
        \tag{$*$}
        12x_1^5x_2^7x_4 - 18x_2^3x_3 + 11x_1^6x_2x_3^3x_4^{23}
    \end{equation}
    is a typical element of $R$. Each $\sigma \in S_4$ gives a permutation of $\set{x_1, \ldots, x_4}$ by defining $\sigma \cdot x_i = x_{\sigma(i)}$. This may be extended to a map from $R$ to $R$ by defining
    \[
    	\sigma \cdot p(x_1, x_2, x_3, x_4) = p(x_{\sigma(1)}, x_{\sigma(2)}, x_{\sigma(3)}, x_{\sigma(4)})
    \]
    for all $p(x_1, x_2, x_3, x_4) \in R$ (i.e., $\sigma$ simply permutes the indices of the variables).
    \begin{subexercises}
        \item Let $p = p(x_1, \ldots, x_4)$ be the polynomial in \eqref{ex2.2.11eq} above, let $\sigma = (1\ 2\ 3\ 4)$ and let $\tau = (1\ 2\ 3)$. Compute $\sigma \cdot p, \tau \cdot (\sigma \cdot p), (\tau \circ \sigma) \cdot p$, and $(\sigma \circ \tau) \cdot p$.
        \item Prove that these definitions give a (left) group action of $S_4$ on $R$.
        \item Exhibit all permutations in $S_4$ that stabilize $x_4$ and prove that they form a subgroup isomorphic to $S_3$.
        \item Exhibit all permutations in $S_4$ that stabilize the element $x_1 + x_2$ and prove that they form an Abelian subgroup of order 4.
        \item Exhibit all permutations in $S_4$ that stabilize the element $x_1x_2 + x_3x_4$ and prove that they form a subgroup isomorphic to the dihedral group of order 8.
        \item Show that the permutations in $S_4$ that stabilize the element $(x_1 + x_2)(x_3 + x_4)$ are exactly the same as those found in part (e).
    \end{subexercises}
    \note{The two polynomials appearing in parts (e) and (f) and the subgroup that stabilizes them will play an important role in the study of roots of quartic equations in Section 14.6.}
\end{exercise}

\begin{solenum}
    \item Note that $\tau \circ \sigma = (1\ 3\ 4\ 2)$ and $\sigma \circ \tau = (1\ 3\ 2\ 4)$. Then
    \begin{align*}
        \sigma \cdot p & = 12x_1x_2^5x_3^7 - 18x_3^3x_4 + 11x_1^{23}x_2^6x_3x_4^3 \\
        \tau \cdot (\sigma \cdot p) & = 12x_1^7x_2x_3^5 - 18x_1^3x_4 + 11x_1x_2^{23}x_3^6x_4^3 \\
        (\tau \circ \sigma) \cdot p & = 12x_1^7x_2x_3^5 - 18x_1^3x_4 + 11x_1x_2^{23}x_3^6x_4^3  \\
        (\sigma \circ \tau) \cdot p & = 12x_1x_3^5x_4^7 - 18x_2x_4^3 + 11x_1^{23}x_2^3x_3^6x_4
    \end{align*}
    \item Note that $1 \cdot p = p$ for any $p \in R$. Let $\sigma, \tau \in S_4$, then
    \begin{align*}
        \sigma \cdot (\tau \cdot p(x_1, x_2, x_3, x_4)) & = \sigma \cdot p(x_{\tau(1)}, x_{\tau(2)}, x_{\tau(3)}, x_{\tau(4)}) \\
        & = p(x_{\sigma(\tau(1))}, x_{\sigma(\tau(2))}, x_{\sigma(\tau(3))}, x_{\sigma(\tau(4))}) \\
        & = (\sigma \circ \tau) \cdot p(x_1, x_2, x_3, x_4)
    \end{align*}
    \item The permutations that stabilize $x_4$ fix 4 and permute $\set{1, 2, 3}$:
    \[
    	G = \set{1, (1\ 2), (1\ 3), (2\ 3), (1\ 2\ 3), (1\ 3\ 2)} \cong S_3.
    \]
    \item The permutations that stabilize $x_1 + x_2$ swap 1 and 2 or fix them:
    \[
    	H = \set{1, (1\ 2), (3\ 4), (1\ 2)(3\ 4)}
    \]
    which is Abelian since every non-identity element has order 2. Moreover, $|H| = 4$.
    \item The permutations that stabilize $x_1x_2 + x_3x_4$ swap 1 and 2 or fix them, swap 3 and 4 or fix them, swap the pairs $(1\ 2)$ and $(3\ 4)$, or swap all four elements:
    \[
    	K = \set{1, (1\ 2), (3\ 4), (1\ 2)(3\ 4), (1\ 3)(2\ 4), (1\ 4)(2\ 3), (1\ 3\ 2\ 4), (1\ 4\ 2\ 3)}.
    \]
    It is straightforward to see that $K$ is generated by $(1\ 2)$ and $(1\ 3\ 2\ 4)$. Let $\phi : D_8 \to K$ be defined by
    \[
    	\phi(r) = (1\ 3\ 2\ 4), \quad \phi(s) = (1\ 2)
    \]
    Since $\phi(r)^4 = \phi(s)^2 = 1$, and 
    \[
    	\phi(s) \phi(r) = (1\ 2)(1\ 3\ 2\ 4) = (1\ 3)(2\ 4) = (1\ 4\ 2\ 3)(1\ 2) = \phi(r)\inv \phi(s),
    \]
	then $\phi$ is a homomorphism. Since $\phi$ is clearly surjective and $|D_8| = |K| = 8$, then $\phi$ is an isomorphism, so that $K \cong D_8$.
    \item The permutations that stabilize $(x_1 + x_2)(x_3 + x_4)$ must either swap 1 and 2 or fix them, and either swap 3 and 4 or fix them. Additionally, it may swap the pairs $(1\ 2)$ and $(3\ 4)$. These are exactly the same permutations found in part (e). \qh
\end{solenum}

\begin{exercise}[2.2.13]
    Let $n$ be a positive integer and let $R$ be the set of all polynomials with integer coefficients in the independent variables $x_1, x_2, \ldots, x_n$, i.e., the members of $R$ are finite sums of elements of the form $ax_1^{r_1}x_2^{r_2} \cdots x_n^{r_n}$ where $a$ is any integer and $r_1, \ldots, r_n$ are nonnegative integers. For each $\sigma \in S_n$ define a map
    \[
    	\sigma : R \to R \quad \text{by} \quad \sigma \cdot p(x_1, x_2, \ldots, x_n) = p(x_{\sigma(1)}, x_{\sigma(2)}, \ldots, x_{\sigma(n)})
    \]
    Prove that this defines a (left) group action of $S_n$ on $R$.
\end{exercise}

\begin{sol}
    Since $\id(i) = i$ for all $1 \leq i \leq n$, then $\id \cdot p = p$ for any $p \in R$. Let $\sigma, \tau \in S_n$, then
    \begin{align*}
        \sigma \cdot (\tau \cdot p(x_1, x_2, \ldots, x_n)) & = \sigma \cdot p(x_{\tau(1)}, x_{\tau(2)}, \ldots, x_{\tau(n)}) \\
        & = p(x_{\sigma(\tau(1))}, x_{\sigma(\tau(2))}, \ldots, x_{\sigma(\tau(n))}) \\
        & = (\sigma \circ \tau) \cdot p(x_1, x_2, \ldots, x_n)
    \end{align*}
    Hence, this defines a (left) group action of $S_n$ on $R$.
\end{sol}

\begin{exercise}[2.2.14]
    Let $H(F)$ be the Heisenberg group over the field $F$ introduced in \exref{1.4.11}. Determine which matrices lie in the center of $H(F)$ and prove that $Z(H(F))$ is isomorphic to the additive group $F$.
\end{exercise}

\begin{sol}
    Let $X, Y \in H(F)$ be written as
    \[
    	X = (a, b, c) \quad \text{and} \quad Y = (d, e, f)
    \]
    where $X \in Z(H(F))$. Then $XY = YX$, implying that
    \[
    	XY = (a + d, b + e + af, c + f) = (d + a, e + b + dc, f + c) = YX
    \]
    so that $af + b + e = dc + e + b$. Simplifying, we have $af = dc$. Since this must hold for all $d, f \in F$, then $a = 0$ and $c = 0$. It follows that
    \[
    	Z(H(F)) = \set{(0, x, 0) \mid x \in F}
    \]
    Moreover, the mapping $\phi : F \to Z(H(F))$ defined by
    \[
    	\phi(x) = (0, x, 0)
    \]
    is clearly surjective. If $\phi(x) = (0, 0, 0)$, then $x = 0$, so that $\phi$ is injective. Finally, for any $x, y \in F$, we have
    \[
    	\phi(x + y) = (0, x + y, 0) = (0, x, 0)(0, y, 0) = \phi(x) \phi(y).
    \]
    Hence, $\phi$ is an isomorphism, so that $Z(H(F)) \cong F$.
\end{sol}